\chapter{Systeempalet}

\section{Systeempalet}
\label{appendix:colourtable}

De volgende tabel beschrijft het systeempalet zoals het standaard is vastgelegd. U kunt deze paletkleuren wijzigen met het commando {\bf PALETTE COLOR} en terugzetten op hun standaardwaarden met het commando {\bf PALETTE RESTORE}.

Paletindexen worden gebruikt als waarden in de speciale array {\bf C@\&()} en als argumenten voor BASIC-commando's zoals {\bf BACKGROUND}, {\bf BORDER}, {\bf COLOR}, {\bf FOREGROUND}, {\bf HIGHLIGHT}, {\bf PEN} en {\bf SCREEN CLR}.

\definecolor{m65black}{RGB}{0, 0, 0}\definecolor{m65white}{RGB}{255, 255, 255}\definecolor{m65red}{RGB}{255, 0, 0}\definecolor{m65cyan}{RGB}{0, 255, 255}\definecolor{m65purple}{RGB}{255, 0, 255}\definecolor{m65green}{RGB}{0, 255, 0}\definecolor{m65blue}{RGB}{0, 0, 255}\definecolor{m65yellow}{RGB}{255, 255, 0}\definecolor{m65orange}{RGB}{255, 111, 0}\definecolor{m65brown}{RGB}{175, 79, 0}\definecolor{m65lightred}{RGB}{255, 127, 127}\definecolor{m65darkgrey}{RGB}{95, 95, 95}\definecolor{m65mediumgrey}{RGB}{143, 143, 143}\definecolor{m65lightgreen}{RGB}{159, 255, 159}\definecolor{m65lightblue}{RGB}{159, 159, 255}\definecolor{m65lightgrey}{RGB}{191, 191, 191}\definecolor{m65gurumeditation}{RGB}{239, 0, 0}\definecolor{m65rambutan}{RGB}{255, 95, 0}\definecolor{m65carrot}{RGB}{255, 191, 0}\definecolor{m65lemontart}{RGB}{239, 239, 0}\definecolor{m65pandan}{RGB}{127, 255, 0}\definecolor{m65seasickgreen}{RGB}{111, 239, 111}\definecolor{m65soylentgreen}{RGB}{0, 239, 61}\definecolor{m65slimergreen}{RGB}{0, 255, 159}\definecolor{m65theothercyan}{RGB}{0, 223, 223}\definecolor{m65seasky}{RGB}{0, 159, 255}\definecolor{m65smurfblue}{RGB}{0, 63, 255}\definecolor{m65screenofdeath}{RGB}{0, 0, 239}\definecolor{m65plumsauce}{RGB}{127, 0, 255}\definecolor{m65sourgrape}{RGB}{207, 0, 255}\definecolor{m65bubblegum}{RGB}{255, 0, 191}\definecolor{m65hottamales}{RGB}{255, 61, 111}

\begin{center}
    {\setlength{\tabcolsep}{1mm}
    \begin{tabular}{*{4}{|R{1.2cm}}|l|}
    \hline
    {\bf Index}  & {\bf Rood} & {\bf Groen} & {\bf Blauw} & {\bf Kleur} \\
    \hline
      {\bf 0} &    0  &   0   &  0   & \tikz[scale=0.3] \draw[fill=m65black] (0,0) rectangle (2,1); Zwart \\      {\bf 1} &   15  &  15   & 15   & \tikz[scale=0.3] \draw[fill=m65white] (0,0) rectangle (2,1); Wit \\      {\bf 2} &   15  &   0   &  0   & \tikz[scale=0.3] \draw[fill=m65red] (0,0) rectangle (2,1); Rood   \\      {\bf 3} &    0  &  15   & 15   & \tikz[scale=0.3] \draw[fill=m65cyan] (0,0) rectangle (2,1); Cyaan  \\      {\bf 4} &   15  &   0   & 15   & \tikz[scale=0.3] \draw[fill=m65purple] (0,0) rectangle (2,1); Paars\\      {\bf 5} &    0  &  15   &  0   & \tikz[scale=0.3] \draw[fill=m65green] (0,0) rectangle (2,1); Groen \\      {\bf 6} &    0  &   0   & 15   & \tikz[scale=0.3] \draw[fill=m65blue] (0,0) rectangle (2,1); Blauw  \\      {\bf 7} &   15  &  15   &  0   & \tikz[scale=0.3] \draw[fill=m65yellow] (0,0) rectangle (2,1); Geel\\      {\bf 8} &   15  &   6   &  0   & \tikz[scale=0.3] \draw[fill=m65orange] (0,0) rectangle (2,1); Oranje\\      {\bf 9} &   10  &   4   &  0   & \tikz[scale=0.3] \draw[fill=m65brown] (0,0) rectangle (2,1); Bruin \\      {\bf 10} &   15  &   7   &  7   & \tikz[scale=0.3] \draw[fill=m65lightred] (0,0) rectangle (2,1); Lichtrood (roze)  \\      {\bf 11} &    5  &   5   &  5   & \tikz[scale=0.3] \draw[fill=m65darkgrey] (0,0) rectangle (2,1); Donkergrijs\\      {\bf 12} &    8  &   8   &  8   & \tikz[scale=0.3] \draw[fill=m65mediumgrey] (0,0) rectangle (2,1); Middengrijs\\      {\bf 13} &    9  &  15   &  9   & \tikz[scale=0.3] \draw[fill=m65lightgreen] (0,0) rectangle (2,1); Lichtgroen \\      {\bf 14} &    9  &   9   & 15   & \tikz[scale=0.3] \draw[fill=m65lightblue] (0,0) rectangle (2,1); Lichtblauw\\      {\bf 15} &   11  &  11   & 11   & \tikz[scale=0.3] \draw[fill=m65lightgrey] (0,0) rectangle (2,1); Lichtgrijs\\
    \hline
      {\bf 16} &   14  &   0   &  0   & \tikz[scale=0.3] \draw[fill=m65gurumeditation] (0,0) rectangle (2,1); Guru-meditatie\\      {\bf 17} &   15  &   5   &  0   & \tikz[scale=0.3] \draw[fill=m65rambutan] (0,0) rectangle (2,1); Rambtoetan\\      {\bf 18} &   15  &  11   &  0   & \tikz[scale=0.3] \draw[fill=m65carrot] (0,0) rectangle (2,1); Wortel\\      {\bf 19} &   14  &  14   &  0   & \tikz[scale=0.3] \draw[fill=m65lemontart] (0,0) rectangle (2,1); Citroentaart\\      {\bf 20} &    7  &  15   &  0   & \tikz[scale=0.3] \draw[fill=m65pandan] (0,0) rectangle (2,1); Pandangroen\\      {\bf 21} &    6  &  14   &  6   & \tikz[scale=0.3] \draw[fill=m65seasickgreen] (0,0) rectangle (2,1); Zeeziek groen\\      {\bf 22} &    0  &  14   &  3   & \tikz[scale=0.3] \draw[fill=m65soylentgreen] (0,0) rectangle (2,1); Soylent Groen\\      {\bf 23} &    0  &  15   &  9   & \tikz[scale=0.3] \draw[fill=m65slimergreen] (0,0) rectangle (2,1); Slijmgroen\\      {\bf 24} &    0  &  13   &  13  & \tikz[scale=0.3] \draw[fill=m65theothercyan] (0,0) rectangle (2,1); Het andere cyaan\\      {\bf 25} &    0  &   9   &  15  & \tikz[scale=0.3] \draw[fill=m65seasky] (0,0) rectangle (2,1); Zeehemel\\      {\bf 26} &    0  &   3   &  15  & \tikz[scale=0.3] \draw[fill=m65smurfblue] (0,0) rectangle (2,1); Smurfenblauw\\      {\bf 27} &    0  &   0   &  14  & \tikz[scale=0.3] \draw[fill=m65screenofdeath] (0,0) rectangle (2,1); Scherm des Doods\\      {\bf 28} &    7  &   0   &  15  & \tikz[scale=0.3] \draw[fill=m65plumsauce] (0,0) rectangle (2,1); Pruimensaus\\      {\bf 29} &   12  &   0   &  15  & \tikz[scale=0.3] \draw[fill=m65sourgrape] (0,0) rectangle (2,1); Zure druif\\      {\bf 30} &   15  &   0   &  11  & \tikz[scale=0.3] \draw[fill=m65bubblegum] (0,0) rectangle (2,1); Kauwgom\\      {\bf 31} &   15  &   3   &   6  & \tikz[scale=0.3] \draw[fill=m65hottamales] (0,0) rectangle (2,1); Hete tamales\\
    \hline
    \end{tabular}
    }
    \end{center}

U kunt deze 32 kleuren in de schermterminal gebruiken door PETSCII-Controlcodes te typen of weer te geven. De tekenstatus van de cursor is ofwel één van de eerste 16 kleuren ofwel één van de tweede 16 kleuren, en deze status kan met Controlcodes worden gewijzigd. Elk van de 16 Controlcodes kiest de kleur van de cursor binnen die kleurenset.

Om kleurenset A te kiezen, drukt u op \specialkey{CTRL} + \megakey{D} of geeft u PETSCII-code 4 weer.

Om kleurenset B te kiezen, drukt u op \specialkey{CTRL} + \megakey{A} of geeft u PETSCII-code 1 weer.

\begin{center}
  {\setlength{\tabcolsep}{1mm}\renewcommand{\arraystretch}{1.5}
  \begin{tabular}{|l|R{1.3cm}|l|l|}
  \hline
  {\bf Toets} & {\bf PETSCII} & {\bf Kleurenset A} & {\bf Kleurenset B} \\
  \hline
  \specialkey{CTRL} + \megakey{1} & 144 & \tikz[scale=0.3] \draw[fill=m65black] (0,0) rectangle (2,1); Zwart & \tikz[scale=0.3] \draw[fill=m65gurumeditation] (0,0) rectangle (2,1); Guru-meditatie\\
  \hline
  \specialkey{CTRL} + \megakey{2} & 5 & \tikz[scale=0.3] \draw[fill=m65white] (0,0) rectangle (2,1); Wit & \tikz[scale=0.3] \draw[fill=m65rambutan] (0,0) rectangle (2,1); Rambtoetan\\
  \hline
  \specialkey{CTRL} + \megakey{3} & 28 & \tikz[scale=0.3] \draw[fill=m65red] (0,0) rectangle (2,1); Rood & \tikz[scale=0.3] \draw[fill=m65carrot] (0,0) rectangle (2,1); Wortel\\
  \hline
  \specialkey{CTRL} + \megakey{4} & 159 & \tikz[scale=0.3] \draw[fill=m65cyan] (0,0) rectangle (2,1); Cyaan & \tikz[scale=0.3] \draw[fill=m65lemontart] (0,0) rectangle (2,1); Citroentaart\\
  \hline
  \specialkey{CTRL} + \megakey{5} & 156 & \tikz[scale=0.3] \draw[fill=m65purple] (0,0) rectangle (2,1); Paars & \tikz[scale=0.3] \draw[fill=m65pandan] (0,0) rectangle (2,1); Pandangroen\\
  \hline
  \specialkey{CTRL} + \megakey{6} & 30 & \tikz[scale=0.3] \draw[fill=m65green] (0,0) rectangle (2,1); Groen & \tikz[scale=0.3] \draw[fill=m65seasickgreen] (0,0) rectangle (2,1); Zeeziek groen\\
  \hline
  \specialkey{CTRL} + \megakey{7} & 31 & \tikz[scale=0.3] \draw[fill=m65blue] (0,0) rectangle (2,1); Blauw & \tikz[scale=0.3] \draw[fill=m65soylentgreen] (0,0) rectangle (2,1); Soylent Groen\\
  \hline
  \specialkey{CTRL} + \megakey{8} & 158 & \tikz[scale=0.3] \draw[fill=m65yellow] (0,0) rectangle (2,1); Geel & \tikz[scale=0.3] \draw[fill=m65slimergreen] (0,0) rectangle (2,1); Slijmgroen\\
  \hline
  \megasymbolkey + \megakey{1} & 129 & \tikz[scale=0.3] \draw[fill=m65orange] (0,0) rectangle (2,1); Oranje & \tikz[scale=0.3] \draw[fill=m65theothercyan] (0,0) rectangle (2,1); Het andere cyaan\\
  \hline
  \megasymbolkey + \megakey{2} & 149 & \tikz[scale=0.3] \draw[fill=m65brown] (0,0) rectangle (2,1); Bruin & \tikz[scale=0.3] \draw[fill=m65seasky] (0,0) rectangle (2,1); Zeehemel\\
  \hline
  \megasymbolkey + \megakey{3} & 150 & \tikz[scale=0.3] \draw[fill=m65lightred] (0,0) rectangle (2,1); Lichtrood (roze) & \tikz[scale=0.3] \draw[fill=m65smurfblue] (0,0) rectangle (2,1); Smurfenblauw\\
  \hline
  \megasymbolkey + \megakey{4} & 151 & \tikz[scale=0.3] \draw[fill=m65darkgrey] (0,0) rectangle (2,1); Donkergrijs & \tikz[scale=0.3] \draw[fill=m65screenofdeath] (0,0) rectangle (2,1); Scherm des Doods\\
  \hline
  \megasymbolkey + \megakey{5} & 152 & \tikz[scale=0.3] \draw[fill=m65mediumgrey] (0,0) rectangle (2,1); Middengrijs & \tikz[scale=0.3] \draw[fill=m65plumsauce] (0,0) rectangle (2,1); Pruimensaus\\
  \hline
  \megasymbolkey + \megakey{6} & 153 & \tikz[scale=0.3] \draw[fill=m65lightgreen] (0,0) rectangle (2,1); Lichtgroen & \tikz[scale=0.3] \draw[fill=m65sourgrape] (0,0) rectangle (2,1); Zure druif\\
  \hline
  \megasymbolkey + \megakey{7} & 154 & \tikz[scale=0.3] \draw[fill=m65lightblue] (0,0) rectangle (2,1); Lichtblauw & \tikz[scale=0.3] \draw[fill=m65bubblegum] (0,0) rectangle (2,1); Kauwgom\\
  \hline
  \megasymbolkey + \megakey{8} & 155 & \tikz[scale=0.3] \draw[fill=m65lightgrey] (0,0) rectangle (2,1); Lichtgrijs & \tikz[scale=0.3] \draw[fill=m65hottamales] (0,0) rectangle (2,1); Hete tamales\\
  \hline
\end{tabular}
  }
\end{center}
